%═══════════════════════════════════════════════════════════════════════════════
\chapter{La ecuación de Fokker-Planck y la reversión temporal de SDEs}
\label{ch:cap11}
%═══════════════════════════════════════════════════════════════════════════════

El Capítulo~\ref{ch:cap10} terminó enunciando la ecuación de Fokker-Planck sin derivarla y
verificando que la distribución de Gibbs del Capítulo~\ref{ch:cap5} era estacionaria de la
SDE de Langevin sin demostración rigurosa. Este capítulo salda esas deudas y
construye sobre ellas el resultado que organiza los modelos generativos basados
en difusión: el teorema de Anderson sobre la reversión temporal de SDEs.

Hay una observación de fondo, relevante para la física de reactores y para
cualquier sistema estocástico con correlaciones espaciales: la difusión isótropa
$\bm{\Sigma} = g(t)\bm{I}$ que usan los modelos de difusión modernos no es la
única opción posible, ni necesariamente la más eficiente. El teorema de Anderson
se establece primero en su forma \textbf{más general} —con difusión $\bm{\Sigma}(\bm{x},t)$
dependiente del estado— y el caso isótropo aparece como corolario. El término
$\nabla\cdot\bm{D}$ que distingue al caso general tiene interpretación física
precisa: es la fuerza de difusio\-fo\-ré\-ti\-ca que emerge cuando el tensor de
difusión varía en el espacio, análoga al término de deriva en la ecuación de
transporte de neutrones en medios heterogéneos.

El capítulo cierra con el puente de Schrödinger: el proceso reverso de Anderson
no es solo una SDE con marginales correctas, sino la única que minimiza la KL
entre leyes de procesos. En el límite de difusión cero, ese puente converge
al mapa de transporte óptimo de Monge-Kantorovich, conectando la teoría de
difusión con el \textit{flow matching} de los modelos generativos más recientes.

%───────────────────────────────────────────────────────────────────────────────
\section{La ecuación de Fokker-Planck: derivación rigurosa}
\label{sec:fp}
\dificultad{3}{3}
%───────────────────────────────────────────────────────────────────────────────

\subsection{Del generador al adjunto}
\label{subsec:fp_derivacion}

Sea $d\bm{X}_t = \bm{\mu}(\bm{X}_t,t)\,dt + \bm{\Sigma}(\bm{X}_t,t)\,d\bm{W}_t$
con $\bm{W}_t \in \mathbb{R}^m$, tensor de difusión $\bm{D} =
\bm{\Sigma}\bm{\Sigma}^\top \in \mathbb{R}^{d\times d}$, y densidad marginal
$p_t(\bm{x})$. El generador infinitesimal actúa sobre $f \in C_c^2(\mathbb{R}^d)$:
\[
  \mathcal{A}f = \sum_i \mu_i\,\partial_{x_i}f
  + \frac{1}{2}\sum_{i,j} D_{ij}\,\partial_{x_i x_j}f.
\]
Por la Proposición~\ref{prop:generador_def}:
$\frac{d}{dt}\mathbb{E}[f(\bm{X}_t)] = \int \mathcal{A}f\,p_t\,d\bm{x}$.
Integrando por partes dos veces y usando que $f$ tiene soporte compacto:
\begin{align*}
  \int \mu_i\,\partial_{x_i}f\cdot p_t &= -\int f\,\partial_{x_i}(\mu_i p_t),\\
  \int D_{ij}\,\partial_{x_ix_j}f\cdot p_t &= \int f\,\partial_{x_ix_j}(D_{ij}p_t).
\end{align*}

\begin{teorema}{Ecuación de Fokker-Planck}
\label{thm:fokker_planck}
La densidad marginal $p_t$ del proceso de Itô con drift $\bm{\mu}$ y tensor de
difusión $\bm{D} = \bm{\Sigma}\bm{\Sigma}^\top$ satisface:
\begin{equation}
  \frac{\partial p_t}{\partial t}
  = \mathcal{A}^* p_t
  := -\sum_i \partial_{x_i}(\mu_i p_t)
  + \frac{1}{2}\sum_{i,j}\partial_{x_i x_j}(D_{ij}p_t).
  \label{eq:fp}
\end{equation}
\end{teorema}

\subsection{Identidad estructural del término de segundo orden}
\label{subsec:fp_identidad}

Para la demostración del teorema de Anderson será indispensable la siguiente
identidad de expansión del término difusivo. Expandiendo la derivada doble:
\begin{equation}
  \frac{1}{2}\sum_{i,j}\partial_{x_ix_j}(D_{ij}p)
  = \frac{1}{2}\nabla\cdot\!\left[(\nabla\cdot\bm{D})\,p
  + \bm{D}\nabla p\right],
  \label{eq:identidad_difusiva}
\end{equation}
donde $(\nabla\cdot\bm{D})_i = \sum_j\partial_{x_j}D_{ij}$ es la
\textbf{divergencia por filas} del tensor de difusión.

\begin{proof}
Expandiendo $\sum_j\partial_{x_j}(D_{ij}p) = (\nabla\cdot\bm{D})_i p +
\sum_j D_{ij}\partial_{x_j}p$, y tomando $\partial_{x_i}$ y sumando en $i$:
$\sum_{ij}\partial_{x_ix_j}(D_{ij}p) = \sum_i\partial_{x_i}[(\nabla\cdot\bm{D})_i p
+ (\bm{D}\nabla p)_i] = \nabla\cdot[(\nabla\cdot\bm{D})p + \bm{D}\nabla p]$.
$\square$
\end{proof}

La Fokker-Planck en forma de ecuación de continuidad es $\partial_t p_t +
\nabla\cdot\bm{J}_t = 0$ con corriente:
\begin{equation}
  \bm{J}_t = \bm{\mu}\,p_t - \frac{1}{2}[(\nabla\cdot\bm{D})\,p_t + \bm{D}\nabla p_t]
  = p_t\!\left[\bm{\mu} - \frac{1}{2}(\nabla\cdot\bm{D}) - \frac{1}{2}\bm{D}\nabla\log p_t\right].
  \label{eq:corriente}
\end{equation}

\subsection{Distribución estacionaria de la SDE de Langevin}
\label{subsec:gibbs_verificacion}

\begin{proposicion}{La distribución de Gibbs es estacionaria}
\label{prop:gibbs_estacionaria}
Para la SDE $d\bm{\theta}_t = -\nabla\mathcal{L}(\bm{\theta}_t)\,dt +
\sigma\,d\bm{W}_t$ con $\sigma$ constante, $p^* \propto e^{-2\mathcal{L}/\sigma^2}$
satisface $\mathcal{A}^*p^* = 0$.
\end{proposicion}

\begin{proof}
$\nabla p^* = -(2/\sigma^2)\nabla\mathcal{L}\cdot p^*$. Sustituyendo en~\eqref{eq:fp}
con $\bm{D} = \sigma^2\bm{I}$ (para el que $\nabla\cdot\bm{D} = \bm{0}$):
\begin{align*}
  \mathcal{A}^*p^*
  &= \nabla\cdot(\nabla\mathcal{L}\cdot p^*) + \tfrac{\sigma^2}{2}\Delta p^* \\
  &= (\Delta\mathcal{L})p^* - \tfrac{2}{\sigma^2}\|\nabla\mathcal{L}\|^2 p^*
    - (\Delta\mathcal{L})p^* + \tfrac{2}{\sigma^2}\|\nabla\mathcal{L}\|^2 p^* = 0.
    \quad\square
\end{align*}
\end{proof}

\begin{fisicaconexion}{La ecuación de Fokker-Planck en física estadística}
La ecuación de Fokker-Planck~\eqref{eq:fp} tiene una larga historia en física
antes de su adopción en los modelos de difusión de ML. Fokker (1914) y Planck
(1917) la derivaron en el contexto de la difusión browniana. En
\textbf{física estadística de no equilibrio}, la Fokker-Planck es la ecuación
maestra continua: describe la evolución de la distribución de probabilidad de
un sistema estocástico bajo la acción simultánea de deriva determinista y ruido
térmico. La corriente de probabilidad~\eqref{eq:corriente} tiene interpretación
termodinámica directa: $\bm{J} = \bm{0}$ en la distribución estacionaria es el
principio de balance detallado, que equivale a la ausencia de corrientes
circulantes en el equilibrio termodinámico (segunda ley de la termodinámica para
sistemas de Markov).

La Proposición~\ref{prop:gibbs_estacionaria} es el enunciado preciso de la
equipartición estocástica: para la SDE de Langevin con potencial $\mathcal{L}$,
la distribución de Gibbs $p^* \propto e^{-2\mathcal{L}/\sigma^2}$ es la
distribución de equilibrio. El parámetro $\sigma^2/2$ juega el rol de $k_BT$
(temperatura): el ratio $\mathcal{L}/(k_BT)$ determina cuánto pesa cada
configuración en el ensemble estadístico. Esta es la misma distribución de
Gibbs-Boltzmann de la mecánica estadística: el SGD con ruido explora el paisaje
de pérdida exactamente como un sistema físico a temperatura $T_{\mathrm{eff}} =
\sigma^2/2$ explora su espacio de fases. La tasa de disipación de la KL
$\frac{d}{dt}D_{\mathrm{KL}}(p_t\|p^*) = -\frac{\sigma^2}{2}\mathcal{I}(p_t\|p^*)$
(Proposición~\ref{prop:kl_decrece}) tiene la estructura del \textbf{H-teorema
de Boltzmann}: la entropía relativa $H = D_{\mathrm{KL}}$ decrece monótonamente
hasta el equilibrio, con la información de Fisher $\mathcal{I}$ como tasa de
producción de entropía.
\end{fisicaconexion}
\label{sec:balance_detallado}
\dificultad{3}{3}
%───────────────────────────────────────────────────────────────────────────────

\subsection{Balance detallado general}
\label{subsec:bd_general}

La distribución estacionaria satisface $\nabla\cdot\bm{J}^* = 0$. La condición
más fuerte, $\bm{J}^* = \bm{0}$, es el \textbf{balance detallado}. De~\eqref{eq:corriente}:
\begin{equation}
  \bm{J}^* = \bm{0}
  \;\Leftrightarrow\;
  \bm{\mu}(\bm{x}) = \frac{1}{2}(\nabla\cdot\bm{D})(\bm{x})
  + \frac{1}{2}\bm{D}(\bm{x})\nabla\log p^*(\bm{x}).
  \label{eq:balance_general}
\end{equation}
Para $\bm{D} = \sigma^2\bm{I}$ esto se reduce a $\bm{\mu} = \frac{\sigma^2}{2}
\nabla\log p^*$, la condición de la Sección~\ref{subsec:gibbs_verificacion}. Para $\bm{D}$
dependiente del estado aparece el término $\frac{1}{2}\nabla\cdot\bm{D}$:
la fuerza de \textbf{difusioforesis} que arrastra partículas hacia las regiones
de mayor difusividad, independientemente del potencial. Este término tiene
importancia directa en sistemas de transporte en medios heterogéneos.

\begin{proposicion}{Decrecimiento de la KL}
\label{prop:kl_decrece}
Para un proceso en balance detallado respecto a $p^*$:
$\frac{d}{dt}D_{\mathrm{KL}}(p_t\|p^*) = -\int p_t\,
\bm{J}_t^2/(p_t^2)\,d\bm{x} \leq 0$.
\end{proposicion}

Para difusión isótropa constante, esto equivale a
$\frac{d}{dt}D_{\mathrm{KL}} = -\frac{\sigma^2}{2}\mathcal{I}(p_t\|p^*)$
(demostración idéntica a la Sección anterior).

%───────────────────────────────────────────────────────────────────────────────
\section{El teorema de Anderson: forma general}
\label{sec:anderson}
\dificultad{3}{3}
%───────────────────────────────────────────────────────────────────────────────

\subsection{La pregunta}
\label{subsec:reverso_pregunta}

Sea $\{\bm{X}_t\}_{t\in[0,T]}$ con SDE general $d\bm{X}_t =
\bm{\mu}(\bm{X}_t,t)\,dt + \bm{\Sigma}(\bm{X}_t,t)\,d\bm{W}_t$ y densidades
marginales $p_t > 0$. El proceso reverso $\tilde{\bm{X}}_t = \bm{X}_{T-t}$
tiene densidad marginal $\tilde{p}_t = p_{T-t}$. ¿Satisface $\tilde{\bm{X}}_t$
una SDE? ¿Cuál?

\subsection{Enunciado: caso general con difusión dependiente del estado}
\label{subsec:anderson_general}

\begin{teorema}{Reversión temporal de SDEs -- forma general (Anderson, 1982)}
\label{thm:anderson_general}
Sea $d\bm{X}_t = \bm{\mu}(\bm{X}_t,t)\,dt + \bm{\Sigma}(\bm{X}_t,t)\,d\bm{W}_t$
con $\bm{D} = \bm{\Sigma}\bm{\Sigma}^\top$ y $p_t \in C^{1,2}$ estrictamente
positiva. El proceso reverso $\tilde{\bm{X}}_t = \bm{X}_{T-t}$ satisface:

\begin{equation}
\boxed{
\begin{aligned}
d\tilde{\bm{X}}_t
= \Bigl[ & -\bm{\mu}(\tilde{\bm{X}}_t,T{-}t)
+ (\nabla_{\bm{x}}\cdot\bm{D})(\tilde{\bm{X}}_t,T{-}t) \\
& + \bm{D}(\tilde{\bm{X}}_t,T{-}t)\,\nabla_{\bm{x}}\log p_{T-t}(\tilde{\bm{X}}_t)
\Bigr]dt \\
& + \bm{\Sigma}(\tilde{\bm{X}}_t,T{-}t)\,d\tilde{\bm{W}}_t,
\end{aligned}
}
\label{eq:reverse_general}
\end{equation}

donde $(\nabla\cdot\bm{D})_i = \sum_j\partial_{x_j}D_{ij}$ y $\tilde{\bm{W}}_t$
es un movimiento browniano bajo la ley del proceso reverso.
\end{teorema}

El drift del proceso reverso tiene tres contribuciones:
\begin{enumerate}
  \item $-\bm{\mu}$: el drift forward negado (reversión naive).
  \item $\nabla\cdot\bm{D}$: corrección por heterogeneidad espacial de la
        difusión. Cero si $\bm{D}$ no depende del estado.
  \item $\bm{D}\nabla\log p_{T-t}$: corrección por la estructura de la
        distribución marginal. Es el término de score ponderado por la difusión.
\end{enumerate}

\subsection{Demostración: Ruta A — Verificación directa vía Fokker-Planck}
\label{subsec:anderson_ruta_a}

La estrategia es proponer la SDE~\eqref{eq:reverse_general}, escribir su
Fokker-Planck para $\tilde{p}_t$, y verificar que $\tilde{p}_t = p_{T-t}$
la satisface. La verificación es un cálculo algebraico puro que no requiere
teoría de medidas avanzada.

\medskip
\noindent\textbf{Paso 1: Fokker-Planck del proceso reverso propuesto.}

El drift de~\eqref{eq:reverse_general} es
$\tilde{\bm{\mu}}(\bm{x},t) = -\bm{\mu}(\bm{x},T{-}t) + (\nabla\cdot\bm{D})(\bm{x},T{-}t)
+ \bm{D}(\bm{x},T{-}t)\nabla\log p_{T-t}(\bm{x})$
y la difusión es $\bm{\Sigma}(\bm{x},T{-}t)$. La Fokker-Planck para la densidad
$\tilde{p}_t$ del proceso reverso es:
\begin{equation}
  \partial_t\tilde{p}_t
  = -\nabla\cdot(\tilde{\bm{\mu}}\,\tilde{p}_t)
  + \frac{1}{2}\sum_{i,j}\partial_{x_ix_j}(D_{ij}\,\tilde{p}_t).
  \label{eq:fp_reverso}
\end{equation}

\medskip
\noindent\textbf{Paso 2: lado izquierdo con la candidata $\tilde{p}_t = p_{T-t}$.}

Por la regla de la cadena con $\tau = T-t$:
\[
  \partial_t\tilde{p}_t(\bm{x}) = -(\partial_\tau p_\tau)(\bm{x})\big|_{\tau=T-t}.
\]
Como $p_\tau$ satisface la Fokker-Planck forward~\eqref{eq:fp}, usando la
identidad~\eqref{eq:identidad_difusiva}:
\begin{equation}
  \partial_t\tilde{p}_t
  = \nabla\cdot(\bm{\mu}\,p_{T-t})
  - \frac{1}{2}\nabla\cdot\!\left[(\nabla\cdot\bm{D})\,p_{T-t}
    + \bm{D}\nabla p_{T-t}\right].
  \label{eq:lhs_anderson}
\end{equation}
Todos los coeficientes se evalúan en $(\bm{x},T{-}t)$.

\medskip
\noindent\textbf{Paso 3: lado derecho de la Fokker-Planck reversa.}

Sustituyendo $\tilde{p}_t = p_{T-t}$ y $\tilde{\bm{\mu}} = -\bm{\mu} +
(\nabla\cdot\bm{D}) + \bm{D}\nabla\log p_{T-t}$ en~\eqref{eq:fp_reverso}:

\textit{Término de drift:}
\begin{align*}
  -\nabla\cdot(\tilde{\bm{\mu}}\,p_{T-t})
  &= \nabla\cdot(\bm{\mu}\,p_{T-t})
    - \nabla\cdot\!\left[(\nabla\cdot\bm{D})\,p_{T-t}\right]
    - \nabla\cdot(\bm{D}\,p_{T-t}\,\nabla\log p_{T-t}).
\end{align*}
Usando $p_{T-t}\nabla\log p_{T-t} = \nabla p_{T-t}$:
\[
  \nabla\cdot(\bm{D}\,p_{T-t}\,\nabla\log p_{T-t})
  = \nabla\cdot(\bm{D}\nabla p_{T-t}).
\]

\textit{Término difusivo} (con la identidad~\eqref{eq:identidad_difusiva}):
\[
  \frac{1}{2}\sum_{ij}\partial_{x_ix_j}(D_{ij}p_{T-t})
  = \frac{1}{2}\nabla\cdot\!\left[(\nabla\cdot\bm{D})\,p_{T-t}
    + \bm{D}\nabla p_{T-t}\right].
\]

\textit{Suma del lado derecho:}
\begin{align}
  \text{RHS}
  &= \nabla\cdot(\bm{\mu}\,p_{T-t})
    - \nabla\cdot\!\left[(\nabla\cdot\bm{D})\,p_{T-t}\right]
    - \nabla\cdot(\bm{D}\nabla p_{T-t}) \notag\\
  &\quad + \frac{1}{2}\nabla\cdot\!\left[(\nabla\cdot\bm{D})\,p_{T-t}
    + \bm{D}\nabla p_{T-t}\right] \notag\\
  &= \nabla\cdot(\bm{\mu}\,p_{T-t})
    - \frac{1}{2}\nabla\cdot\!\left[(\nabla\cdot\bm{D})\,p_{T-t}
    + \bm{D}\nabla p_{T-t}\right].
  \label{eq:rhs_anderson}
\end{align}

\medskip
\noindent\textbf{Paso 4: comparación.}

Comparando~\eqref{eq:lhs_anderson} y~\eqref{eq:rhs_anderson}: LHS $=$ RHS.
La candidata $\tilde{p}_t = p_{T-t}$ satisface la Fokker-Planck del proceso
reverso propuesto para cualquier tensor $\bm{D}(\bm{x},t)$. La condición
inicial $\tilde{p}_0 = p_T$ es correcta por construcción. Por unicidad de
soluciones de la Fokker-Planck, la SDE~\eqref{eq:reverse_general} es la SDE
del proceso reverso. $\square$

\begin{observacion}{El rol de la identidad $\nabla\cdot(p\nabla\log p) = \nabla\cdot(\nabla p)$}
La pieza algebraica central de la demostración es que el término de score del
drift reverso, al multiplicarse por $p$, produce $\bm{D}\nabla p$, que junto
con el término difusivo de segundo orden de la Fokker-Planck produce exactamente
la mitad de $\nabla\cdot[(\nabla\cdot\bm{D})p + \bm{D}\nabla p]$. La otra
mitad cancela con $-\nabla\cdot[(\nabla\cdot\bm{D})p]$ del término de drift.
El score $\nabla\log p$ codifica la información de curvatura de la distribución
en un campo vectorial de primer orden, lo que permite que la Fokker-Planck de
segundo orden del reverso coincida exactamente con la del forward invertido en
tiempo.
\end{observacion}

\subsection{Demostración: Ruta B — Cambio de medida vía Girsanov}
\label{subsec:anderson_ruta_b}

La Ruta A es algebraicamente directa pero no ilumina la estructura probabilística.
La Ruta B revela por qué el score aparece: es la derivada de Radon-Nikodym entre
la ley del proceso forward y la ley del browniano de referencia.

\begin{definicion}{Exponencial de Doléans-Dade}
\label{def:doleans}
Para $\bm{u}_t \in \mathbb{R}^m$ adaptado, la \textbf{exponencial de
Doléans-Dade} es:
\[
  \mathcal{E}(\bm{u})_t
  = \exp\!\left(\int_0^t \bm{u}_s^\top d\bm{W}_s
    - \frac{1}{2}\int_0^t \|\bm{u}_s\|^2\,ds\right),
\]
que satisface $d\mathcal{E}(\bm{u})_t = \mathcal{E}(\bm{u})_t\,\bm{u}_t^\top
d\bm{W}_t$: es una martingala local.
\end{definicion}

La \textbf{condición de Novikov} $\mathbb{E}\!\left[e^{\frac{1}{2}\int_0^T
\|\bm{u}_t\|^2\,dt}\right] < \infty$ garantiza que $\mathcal{E}(\bm{u})_T$
es una martingala con $\mathbb{E}[\mathcal{E}(\bm{u})_T] = 1$.

\begin{teorema}{Teorema de Girsanov}
\label{thm:girsanov}
Bajo la condición de Novikov, si se define $\tilde{\mathbb{P}}$ por
$d\tilde{\mathbb{P}}/d\mathbb{P}|_{\mathcal{F}_T} = \mathcal{E}(\bm{u})_T$,
entonces $\tilde{\bm{W}}_t = \bm{W}_t - \int_0^t\bm{u}_s\,ds$ es un
movimiento browniano estándar bajo $\tilde{\mathbb{P}}$.
\end{teorema}

\begin{proof}[Demostración del Teorema~\ref{thm:anderson_general} vía Girsanov]
\textbf{Paso 1: medida del proceso reverso.}
La ley de $\tilde{\bm{X}}_t = \bm{X}_{T-t}$ bajo $\mathbb{P}^{\mathrm{fwd}}$
tiene densidad de transición reversa:
\[
  \tilde{p}(\bm{x},t|\bm{y},s) = p(\bm{y},T-s|\bm{x},T-t)
  \cdot\frac{p_{T-t}(\bm{x})}{p_{T-s}(\bm{y})}, \quad s < t.
\]
Se define la medida $\tilde{\mathbb{P}}$ sobre el proceso reverso mediante
$d\tilde{\mathbb{P}}/d\mathbb{P}^{\mathrm{rev}} = \mathcal{E}(\bm{u})_T$
donde $\mathbb{P}^{\mathrm{rev}}$ es la ley del browniano reverso
$\hat{\bm{W}}_t = \bm{W}_{T-t}-\bm{W}_T$.

\textbf{Paso 2: ecuación de $\tilde{\bm{X}}_t$ bajo $\mathbb{P}$.}
Con $\hat{\bm{W}}_t = \bm{W}_{T-t}-\bm{W}_T$ (que es browniano bajo $\mathbb{P}$)
e invirtiendo el tiempo en la SDE:
\[
  d\tilde{\bm{X}}_t = -\bm{\mu}(\tilde{\bm{X}}_t,T{-}t)\,dt
  - \bm{\Sigma}(\tilde{\bm{X}}_t,T{-}t)\,d\hat{\bm{W}}_t.
\]
Bajo $\mathbb{P}$, el drift efectivo es $-\bm{\mu}$ y la difusión es $-\bm{\Sigma}$
(o $+\bm{\Sigma}$ por simetría del browniano).

\textbf{Paso 3: identificación del proceso $\bm{u}$.}
Para que bajo $\tilde{\mathbb{P}}$ el drift sea el correcto, se necesita
$\bm{\Sigma}\bm{u}_t = (\nabla\cdot\bm{D}) + \bm{D}\nabla\log p_{T-t}$.
Dado que $\bm{D} = \bm{\Sigma}\bm{\Sigma}^\top$, si $\bm{\Sigma}$ es
cuadrada e invertible: $\bm{u}_t = \bm{\Sigma}^\top(\bm{\Sigma}\bm{\Sigma}^\top)^{-1}
[(\nabla\cdot\bm{D}) + \bm{D}\nabla\log p_{T-t}] = \bm{\Sigma}^{-\top}
[(\nabla\cdot\bm{D}) + \bm{D}\nabla\log p_{T-t}]$.

\textbf{Paso 4: verificación de Novikov.}
Para coeficientes Lipschitz y $p_t$ acotada inferior y superiormente en compactos,
$\|\bm{u}_t\|^2 \leq C(1+\|\tilde{\bm{X}}_t\|^2)$. El control de momentos del
proceso garantiza $\mathbb{E}[e^{\frac{1}{2}\int_0^T\|\bm{u}\|^2}] < \infty$.
Para el proceso OU con $p_t$ gaussiana, la condición es verificable explícitamente
ya que $\nabla\log p_t$ es lineal en $\bm{x}$.

\textbf{Paso 5: conclusión.}
Por Girsanov, $\tilde{\bm{W}}_t = \hat{\bm{W}}_t - \int_0^t\bm{u}_s\,ds$ es
browniano bajo $\tilde{\mathbb{P}}$, y la SDE de $\tilde{\bm{X}}_t$ bajo
$\tilde{\mathbb{P}}$ tiene el drift de~\eqref{eq:reverse_general}. $\square$
\end{proof}

\begin{observacion}{Qué aporta cada ruta}
La \textbf{Ruta A} (Fokker-Planck) es elemental y algorítmica: proponer, calcular,
comparar. No requiere más que la regla de la cadena y la identidad~\eqref{eq:identidad_difusiva}.
La \textbf{Ruta B} (Girsanov) revela la estructura probabilística: el score
$\nabla\log p_t$ es proporcional a la derivada de Radon-Nikodym entre la medida
del proceso forward y la del browniano de referencia. El proceso reverso es la
medida de proceso que "corrige" el browniano reverso para que tenga las marginales
correctas. Esta perspectiva es la que conecta directamente con el puente de
Schrödinger de la Sección~\ref{sec:schrodinger}.
\end{observacion}

\subsection{El término \texorpdfstring{$\nabla\cdot\bm{D}$}{nabla·D}: interpretación física}
\label{subsec:divergencia_d}

El término $(\nabla\cdot\bm{D})_i = \sum_j\partial_{x_j}D_{ij}$ es la fuerza
neta que experimenta una partícula browniana en un medio con difusividad variable,
independientemente del potencial. Tiene nombres distintos en distintos campos:

\begin{itemize}
  \item \textbf{Física estadística}: fuerza de difusioforesis o deriva espuria
        (en la convención de Itô).
  \item \textbf{Transporte de neutrones}: término de hetero\-genei\-dad en la
        ecuación de transporte difusivo $\nabla\cdot(D(\bm{x})\nabla\phi)
        = (\nabla D)\cdot(\nabla\phi) + D\Delta\phi$: el primer término es el
        análogo de $(\nabla\cdot\bm{D})p$.
  \item \textbf{Mecánica cuántica}: en la ecuación de Fokker-Planck para la
        densidad de probabilidad en mecánica estocástica de Nelson, este término
        aparece como parte de la velocidad osmótica.
\end{itemize}

Para los modelos de difusión modernos, $\bm{D} = g(t)^2\bm{I}$ no depende del
estado, luego $\nabla\cdot\bm{D} = \bm{0}$: el término desaparece y el reverse
SDE solo tiene el score. Pero en aplicaciones de física de reactores donde la
difusividad $D(\bm{x})$ varía en el espacio (distintos materiales, heterogeneidades),
el término es esencial para que el proceso reverso tenga las marginales correctas.

\begin{fisicaconexion}{Difusioforesis y transporte en medios heterogéneos}
El término $\nabla\cdot\bm{D}$ que aparece en el reverse SDE general tiene
nombre y descripción experimental en física. En un fluido con gradiente de
temperatura o composición, partículas coloidales se mueven desde regiones de
baja difusividad hacia regiones de alta difusividad incluso en ausencia de
gradiente de presión o concentración: este efecto es la
\textbf{difusioforesis} (o termoforesis cuando el gradiente es térmico). La
fuerza efectiva sobre la partícula es proporcional a $\nabla D(\bm{x})$,
exactamente la contribución $\frac{1}{2}\nabla\cdot\bm{D}$ del balance
detallado~\eqref{eq:balance_general}.

En la \textbf{ecuación de difusión de neutrones}, el flujo $\phi(\bm{r})$ en
un reactor heterogéneo satisface $-\nabla\cdot(D(\bm{r})\nabla\phi) + \Sigma_a\phi
= S$, donde $D(\bm{r})$ varía entre el combustible ($D \approx 1$ cm), el
moderador ($D \approx 0.16$ cm para agua) y el reflector. Expandiendo:
$-D(\bm{r})\Delta\phi - (\nabla D)\cdot(\nabla\phi) + \Sigma_a\phi = S$.
El término $(\nabla D)\cdot(\nabla\phi)$ es el análogo macroscópico de
$\nabla\cdot\bm{D}$ en la Fokker-Planck: en interfaces entre materiales, el
flujo neutrónico sufre una discontinuidad de derivada proporcional al salto
en $D$. Ignorar este término (asumir $D$ constante por zonas y no en las
interfaces) produce el error más común en el modelado de reactores heterogéneos.

En \textbf{procesos de difusión en biofísica}, las proteínas difunden en
membranas celulares con difusividad que varía entre el interior y los bordes de
las balsas lipídicas (lipid rafts): $D(\bm{x})$ puede cambiar por un factor
10 en escala de nanómetros. La velocidad de deriva inducida $v_{\mathrm{drift}}
= \frac{1}{2}\partial_x D(x)$ acumula proteínas en los bordes de las balsas,
un mecanismo de señalización que no requiere energía directa. En los
modelos de difusión anisótropa del cerebro (tractografía por tensor de difusión),
el tensor $\bm{D}(\bm{r})$ varía con la orientación de los axones: ignorar
$\nabla\cdot\bm{D}$ en el proceso inverso produciría artefactos sistemáticos en
la reconstrucción de la conectividad cerebral.
\end{fisicaconexion}
\label{subsec:anderson_casos}

\begin{table}[h]
\centering
\begin{tabular}{llll}
\toprule
\textbf{Caso} & $\bm{\Sigma}$ & $\nabla\cdot\bm{D}$ & \textbf{Drift reverso} \\
\midrule
Isótropo escalar & $g(t)\bm{I}$ & $\bm{0}$ & $-\bm{\mu}+g^2\nabla\log p$ \\
Isótropo variable & $g(\bm{x},t)\bm{I}$ & $2g\nabla g$ & $-\bm{\mu}+2g\nabla g+g^2\nabla\log p$ \\
Diagonal & $\mathrm{diag}(\sigma_i(\bm{x},t))$ & $(\partial_{x_j}\sigma_i^2)_i$ & caso general \\
Tensor completo & $\bm{\Sigma}(\bm{x},t)$ & $\sum_j\partial_{x_j}D_{ij}$ & \eqref{eq:reverse_general} \\
\bottomrule
\end{tabular}
\caption{Casos particulares del teorema de Anderson general.}
\label{tab:anderson_casos}
\end{table}

\begin{proposicion}{Caso isótropo escalar —  Corolario del Teorema~\ref{thm:anderson_general}}

\label{cor:anderson_isotropo}
Para $\bm{\Sigma}(\bm{x},t) = g(t)\bm{I}$ y $\bm{\mu}(\bm{x},t) = \bm{f}(\bm{x},t)$:
\begin{equation}
  d\tilde{\bm{X}}_t
  = \left[-\bm{f}(\tilde{\bm{X}}_t,T{-}t)
  + g(T{-}t)^2\,\nabla_{\bm{x}}\log p_{T-t}(\tilde{\bm{X}}_t)\right]dt
  + g(T{-}t)\,d\tilde{\bm{W}}_t.
  \label{eq:reverse_isotropo}
\end{equation}
\end{proposicion}

\subsection{La ODE de flujo de probabilidad}
\label{subsec:prob_flow_ode}

\begin{proposicion}{ODE de flujo de probabilidad}
\label{prop:prob_flow}
Para el caso isótropo, la ODE determinista:
\begin{equation}
  \frac{d\bm{x}}{dt}
  = \bm{f}(\bm{x},t) - \frac{g(t)^2}{2}\nabla_{\bm{x}}\log p_t(\bm{x})
  \label{eq:prob_flow_ode}
\end{equation}
tiene densidades marginales $p_t(\bm{x})$ idénticas a las del forward process.
\end{proposicion}

\begin{proof}
La ecuación de continuidad de~\eqref{eq:prob_flow_ode} es $\partial_t p =
-\nabla\cdot(p\dot{\bm{x}})$. Sustituyendo:
$\partial_t p = -\nabla\cdot(p\bm{f}) + \frac{g^2}{2}\nabla\cdot(p\nabla\log p)
= -\nabla\cdot(p\bm{f}) + \frac{g^2}{2}\Delta p$,
que es la Fokker-Planck del forward. Misma EDP, misma condición inicial: por
unicidad, misma solución. $\square$
\end{proof}

La ODE de flujo conecta los modelos de difusión con las Neural ODEs del
Capítulo~\ref{ch:cap9}: para muestrear de $p_0$, se integra~\eqref{eq:prob_flow_ode} desde
$t=T$ hasta $t=0$ con condición inicial $\bm{x}(T)\sim\mathcal{N}(\bm{0},\bm{I})$.
Las trayectorias son deterministas pero su distribución marginal es $p_0$.

\subsection{Aplicación al forward process OU}
\label{subsec:anderson_ou}

Para el proceso OU $d\bm{X}_t = -\bm{X}_t\,dt + \sqrt{2}\,d\bm{W}_t$
($\bm{f} = -\bm{x}$, $g = \sqrt{2}$), el Corolario~\ref{cor:anderson_isotropo}
da:
\begin{equation}
  d\tilde{\bm{X}}_t
  = \Big[\tilde{\bm{X}}_t
  + 2\nabla\log p_{T-t}(\tilde{\bm{X}}_t)\Big]dt
  + \sqrt{2}\,d\tilde{\bm{W}}_t.
  \label{eq:reverse_ou}
\end{equation}
El único elemento desconocido es el score $\nabla\log p_{T-t}$: varía de
$-\tilde{\bm{X}}_t$ (gaussiano, simple) en $t=0$ hasta la geometría completa
de $p_0$ en $t=T$.

%───────────────────────────────────────────────────────────────────────────────
\section{Score matching: estimación de la función de score}
\label{sec:score_matching}
\dificultad{3}{2}
%───────────────────────────────────────────────────────────────────────────────

\subsection{La función de score}
\label{subsec:score_def}

\begin{definicion}{Función de score}
La \textbf{función de score} de una distribución con densidad $p$ es
$\bm{s}(\bm{x}) = \nabla_{\bm{x}}\log p(\bm{x})$.
\end{definicion}

El score apunta hacia las regiones de mayor probabilidad. Para $p =
\mathcal{N}(\bm{\mu},\bm{\Sigma})$: $\bm{s}(\bm{x}) = -\bm{\Sigma}^{-1}
(\bm{x}-\bm{\mu})$, que apunta hacia la media. Nótese que el score es
\textbf{no isótropo} para $\bm{\Sigma} \neq \sigma^2\bm{I}$: incluso cuando
el forward process inyecta ruido isótropo, el score de la distribución marginal
$p_t$ hereda la correlación de $p_0$ (ver la discusión de la Sección anterior
sobre por qué la difusión isótropa no destruye la correlación).

\subsection{Score matching implícito (Hyvärinen, 2005)}
\label{subsec:hyvarinen}

El score verdadero $\nabla\log p$ no puede calcularse directamente desde
muestras porque requiere conocer $p$. El score matching implícito elimina esa
dependencia.

\begin{teorema}{Score matching implícito}
\label{thm:hyvarinen}
Para $\bm{s}_{\bm{\phi}}: \mathbb{R}^d\to\mathbb{R}^d$ paramétrico:
\begin{align}
  J_E(\bm{\phi}) &= \tfrac{1}{2}\mathbb{E}_p\|\bm{s}_{\bm{\phi}}-\nabla\log p\|^2
  \quad\text{(no computable)}, \label{eq:jsm_e}\\
  J_I(\bm{\phi}) &= \mathbb{E}_p\!\left[\mathrm{tr}(\nabla_{\bm{x}}\bm{s}_{\bm{\phi}})
  + \tfrac{1}{2}\|\bm{s}_{\bm{\phi}}\|^2\right]
  \quad\text{(computable)}, \label{eq:jsm_i}
\end{align}
satisfacen $J_E(\bm{\phi}) = J_I(\bm{\phi}) + C$ con $C$ independiente de
$\bm{\phi}$. Luego $\argmin J_E = \argmin J_I$.
\end{teorema}

\begin{proof}
Expandiendo $J_E = \frac{1}{2}\mathbb{E}_p\|\bm{s}_{\bm{\phi}}\|^2
- \mathbb{E}_p[\bm{s}_{\bm{\phi}}\cdot\nabla\log p]
+ \frac{1}{2}\mathbb{E}_p\|\nabla\log p\|^2$.
El último término es la constante $C$. Para el término del medio,
usando $\nabla\log p = \nabla p/p$ e integrando por partes (Green en
$\mathbb{R}^d$, términos de frontera nulos para $p$ con decaimiento rápido):
\[
  \int p\,s_{\bm{\phi},i}\,\frac{\partial\log p}{\partial x_i}
  = \int s_{\bm{\phi},i}\,\partial_{x_i}p
  = -\int p\,\partial_{x_i}s_{\bm{\phi},i}.
\]
Sumando en $i$: $\mathbb{E}_p[\bm{s}_{\bm{\phi}}\cdot\nabla\log p]
= -\mathbb{E}_p[\mathrm{tr}(\nabla\bm{s}_{\bm{\phi}})]$.
Luego $J_E = J_I + C$. $\square$
\end{proof}

El costo de $J_I$ es $O(d)$ evaluaciones del jacobiano de $\bm{s}_{\bm{\phi}}$
por muestra. El denoising score matching elimina incluso ese costo.

\subsection{Denoising score matching (Vincent, 2011)}
\label{subsec:dsm}

\begin{teorema}{Denoising score matching}
\label{thm:dsm}
Sea $q_\sigma(\tilde{\bm{x}}|\bm{x}) = \mathcal{N}(\tilde{\bm{x}};\bm{x},\sigma^2\bm{I})$
y $q_\sigma(\tilde{\bm{x}}) = \int q_\sigma(\tilde{\bm{x}}|\bm{x})p(\bm{x})\,d\bm{x}$.
Entonces:
\[
  \argmin_{\bm{\phi}}\,\mathbb{E}_{q_\sigma(\tilde{\bm{x}})}
  \|\bm{s}_{\bm{\phi}}(\tilde{\bm{x}})-\nabla\log q_\sigma(\tilde{\bm{x}})\|^2
  =
  \argmin_{\bm{\phi}}\,\mathbb{E}_{p(\bm{x})\,q_\sigma(\tilde{\bm{x}}|\bm{x})}
  \|\bm{s}_{\bm{\phi}}(\tilde{\bm{x}})-\nabla\log q_\sigma(\tilde{\bm{x}}|\bm{x})\|^2.
\]
El lado derecho no requiere conocer $q_\sigma(\tilde{\bm{x}})$.
\end{teorema}

\begin{proof}
Denotamos $\bm{s} = \bm{s}_{\bm{\phi}}(\tilde{\bm{x}})$, $\bm{s}_0 = \nabla\log
q_\sigma(\tilde{\bm{x}})$ y $\bm{s}_c = \nabla\log q_\sigma(\tilde{\bm{x}}|\bm{x})$.
Se expande:
\[
  \mathbb{E}_{q_\sigma}\|\bm{s}-\bm{s}_0\|^2
  = \mathbb{E}_{q_\sigma}\|\bm{s}\|^2
  - 2\mathbb{E}_{q_\sigma}[\bm{s}\cdot\bm{s}_0]
  + C_1.
\]
El primer término es $\mathbb{E}_{p,q_\sigma|\bm{x}}\|\bm{s}\|^2$.
Para el segundo, usando $q_\sigma(\tilde{\bm{x}})\nabla\log q_\sigma(\tilde{\bm{x}})
= \nabla_{\tilde{\bm{x}}}q_\sigma(\tilde{\bm{x}}) = \int p(\bm{x})
\nabla_{\tilde{\bm{x}}}q_\sigma(\tilde{\bm{x}}|\bm{x})\,d\bm{x}$:
\begin{align*}
  \mathbb{E}_{q_\sigma}[\bm{s}\cdot\bm{s}_0]
  &= \int q_\sigma(\tilde{\bm{x}})\,\bm{s}\cdot\nabla\log q_\sigma(\tilde{\bm{x}})
    \,d\tilde{\bm{x}} \\
  &= \int\!\int p(\bm{x})q_\sigma(\tilde{\bm{x}}|\bm{x})\,\bm{s}
    \cdot\frac{\nabla_{\tilde{\bm{x}}}q_\sigma(\tilde{\bm{x}})}{q_\sigma(\tilde{\bm{x}})}
    \,d\bm{x}\,d\tilde{\bm{x}}.
\end{align*}
Expresando $\nabla\log q_\sigma(\tilde{\bm{x}})$ como promedio condicional
sobre $p(\bm{x}|\tilde{\bm{x}})$ (por la fórmula de Bayes $p(\bm{x}|\tilde{\bm{x}})
\propto p(\bm{x})q_\sigma(\tilde{\bm{x}}|\bm{x})$):
\[
  \nabla\log q_\sigma(\tilde{\bm{x}})
  = \frac{\int p(\bm{x})\nabla_{\tilde{\bm{x}}}q_\sigma(\tilde{\bm{x}}|\bm{x})
    \,d\bm{x}}{q_\sigma(\tilde{\bm{x}})}
  = \mathbb{E}_{p(\bm{x}|\tilde{\bm{x}})}[\nabla\log q_\sigma(\tilde{\bm{x}}|\bm{x})].
\]
Por tanto $\mathbb{E}_{q_\sigma}[\bm{s}\cdot\bm{s}_0]
= \mathbb{E}_{p,q_\sigma|\bm{x}}[\bm{s}\cdot\bm{s}_c]$.
Sustituyendo:
\[
  \mathbb{E}_{q_\sigma}\|\bm{s}-\bm{s}_0\|^2
  = \mathbb{E}_{p,q_\sigma|\bm{x}}\|\bm{s}\|^2
  - 2\mathbb{E}_{p,q_\sigma|\bm{x}}[\bm{s}\cdot\bm{s}_c] + C_1
  = \mathbb{E}_{p,q_\sigma|\bm{x}}\|\bm{s}-\bm{s}_c\|^2
  + C_1 - C_2,
\]
donde $C_2 = \mathbb{E}_{p,q_\sigma|\bm{x}}\|\bm{s}_c\|^2$ no depende de
$\bm{\phi}$. Los $\argmin$ coinciden. $\square$
\end{proof}

\subsubsection*{La pérdida de predicción de ruido (DDPM)}

Para el kernel gaussiano, con $\tilde{\bm{x}} = \bm{x} + \sigma\bm{\varepsilon}$,
$\bm{\varepsilon}\sim\mathcal{N}(\bm{0},\bm{I})$:
$\nabla_{\tilde{\bm{x}}}\log q_\sigma(\tilde{\bm{x}}|\bm{x}) = -\bm{\varepsilon}/\sigma$.
Parametrizando $\bm{s}_{\bm{\phi}}(\tilde{\bm{x}}) =
-\bm{\varepsilon}_{\bm{\phi}}(\tilde{\bm{x}})/\sigma$:
\begin{equation}
  J_{\mathrm{DDPM}}(\bm{\phi})
  = \mathbb{E}_{p(\bm{x}),\bm{\varepsilon}}\bigl[
  \|\bm{\varepsilon}_{\bm{\phi}}(\bm{x}+\sigma\bm{\varepsilon})-\bm{\varepsilon}\|^2
  \bigr].
  \label{eq:ddpm_loss}
\end{equation}
Esta es la función de pérdida de DDPM (Ho et al., 2020): predecir el ruido
añadido. La red $\bm{\varepsilon}_{\bm{\phi}}$ es en la práctica una U-Net o
transformer condicionado al nivel de ruido $\sigma$ (o equivalentemente al
tiempo $t$).

\begin{observacion}{Correlación y el score}
La discusión entre las Secciones~\ref{sec:anderson} y~\ref{sec:score_matching}
revela la separación de responsabilidades del diseño: el ruido isótropo $g(t)\bm{I}$
no destruye la correlación de $p_0$ sino que la traslada a $p_t$ y por ende al
score $\nabla\log p_t$. La red $\bm{\varepsilon}_{\bm{\phi}}$ aprende esa
correlación implícitamente: aunque predice un vector $\bm{\varepsilon}\in
\mathbb{R}^d$, las componentes de su salida son interdependientes a través
de los pesos de la red. El ruido isótropo simplifica $\nabla\cdot\bm{D} =
\bm{0}$ y hace computable la pérdida~\eqref{eq:ddpm_loss}; la geometría
correlacional de los datos queda completamente codificada en la función
aprendida por la red.
\end{observacion}

\begin{fisicaconexion}{La función de score en física estadística y termodinámica}
La función de score $\bm{s}(\bm{x}) = \nabla\log p(\bm{x})$ es un objeto con
significado físico profundo, independiente de su rol en los modelos de difusión.
Para la distribución de Gibbs-Boltzmann $p(\bm{x}) \propto e^{-U(\bm{x})/(k_BT)}$:
\[
  \bm{s}(\bm{x}) = \nabla\log p(\bm{x}) = -\frac{\nabla U(\bm{x})}{k_BT},
\]
que es la \textbf{fuerza termodinámica} normalizada por la temperatura: apunta
hacia los mínimos de energía libre, exactamente como el score apunta hacia las
regiones de alta probabilidad. En la ecuación de Langevin $d\bm{x} = -\nabla U\,dt
+ \sqrt{2k_BT}\,d\bm{W}$, la deriva puede reescribirse como $k_BT\,\bm{s}(\bm{x})$:
el drift es el score multiplicado por la temperatura.

La información de Fisher $\mathcal{I}(p) = \mathbb{E}_p[\|\nabla\log p\|^2]
= \mathbb{E}_p[\|\bm{s}\|^2]$ (la norma cuadrática media del score) aparece
en varios contextos físicos:
\begin{itemize}
  \item En el límite de baja temperatura, $\mathcal{I}(p)$ es proporcional a
        la varianza del gradiente de energía, que mide la rugosidad del paisaje
        de energía.
  \item El corolario de la ecuación de Cramér-Rao establece que $\mathcal{I}^{-1}$
        es la varianza mínima de cualquier estimador insesgado de los parámetros
        de $p$: la información de Fisher es la curvatura de la log-verosimilitud,
        exactamente la matriz de Fisher del Capítulo~\ref{ch:cap7}.
  \item La tasa de decrecimiento de la KL en la Proposición~\ref{prop:kl_decrece},
        $\frac{d}{dt}D_\mathrm{KL} = -\frac{\sigma^2}{2}\mathcal{I}(p_t\|p^*)$,
        es el \textbf{H-teorema cuantitativo}: la información de Fisher de la
        distribución relativa es la tasa de producción de entropía.
\end{itemize}

El score matching implícito (Teorema~\ref{thm:hyvarinen}) puede interpretarse
como la estimación de la fuerza termodinámica sin conocer el potencial $U$
directamente: se aprende $\bm{s} = -\nabla U / k_BT$ desde las posiciones de
las partículas (muestras de $p$), sin resolver la ecuación de Poisson $\Delta U
= \nabla\cdot(\bm{s}\cdot k_BT)$. Esta perspectiva conecta el score matching
con los métodos de física computacional para estimar energías libres a partir
de simulaciones de dinámica molecular.
\end{fisicaconexion}

%───────────────────────────────────────────────────────────────────────────────
\section{El puente de Schrödinger y el transporte óptimo}
\label{sec:schrodinger}
\dificultad{3}{3}
%───────────────────────────────────────────────────────────────────────────────

\subsection{Multiplicidad de los procesos reversos y el principio de selección}
\label{subsec:multiplicidad}

El Teorema~\ref{thm:anderson_general} determina la SDE del proceso reverso
\textit{dado} un forward process específico. Pero hay infinitas SDEs con
marginales $p_T$ en $t=0$ y $p_0$ en $t=T$. ¿En qué sentido es el proceso de
Anderson ``el'' proceso reverso?

El \textbf{problema del puente de Schrödinger} (Schrödinger, 1931; Föllmer,
1988) responde con un principio variacional: entre todos los procesos de Markov
con las marginales correctas, el puente de Schrödinger es el que está más
cerca del proceso de referencia en el sentido de la KL entre leyes de procesos.

\subsection{KL entre leyes de procesos vía Girsanov}
\label{subsec:kl_procesos}

Para dos procesos de Itô con la misma difusión $g(t)$ pero drifts $b(\bm{x},t)$
y $f(\bm{x},t)$, con leyes $\mathbb{P}^b$ y $\mathbb{P}^f$ sobre
$C([0,T];\mathbb{R}^d)$, el teorema de Girsanov establece que la derivada de
Radon-Nikodym es $d\mathbb{P}^b/d\mathbb{P}^f = \mathcal{E}(\bm{u})_T$ con
$\bm{u} = (b-f)/g$. Por tanto:
\begin{equation}
  D_{\mathrm{KL}}(\mathbb{P}^b\|\mathbb{P}^f)
  = \mathbb{E}_{\mathbb{P}^b}\!\left[\log\frac{d\mathbb{P}^b}{d\mathbb{P}^f}\right]
  = \frac{1}{2}\mathbb{E}_{\mathbb{P}^b}\!\left[
    \int_0^T \frac{\|b(\bm{X}_t,t)-f(\bm{X}_t,t)\|^2}{g(t)^2}\,dt
  \right].
  \label{eq:kl_procesos}
\end{equation}
La KL entre leyes de procesos es la energía cuadrática integrada de la
diferencia de drifts a lo largo de las trayectorias.

\begin{definicion}{Puente de Schrödinger}
\label{def:schrodinger}
Sea $\mathbb{R}$ la ley del forward process (proceso de referencia) y
$\Pi(p_T, p_0)$ el conjunto de procesos de Markov con marginal $p_T$ en $t=0$
y marginal $p_0$ en $t=T$. El \textbf{puente de Schrödinger} entre $p_T$ y
$p_0$ con referencia $\mathbb{R}$ es:
\begin{equation}
  \mathbb{P}^* = \argmin_{\mathbb{P}\in\Pi(p_T,p_0)} D_{\mathrm{KL}}(\mathbb{P}\|\mathbb{R}).
  \label{eq:schrodinger_problem}
\end{equation}
\end{definicion}

\begin{proposicion}{El proceso de Anderson es el puente de Schrödinger}
\label{prop:anderson_schrodinger}
Con la ley del forward process OU como referencia $\mathbb{R}$, el proceso
reverso de Anderson $\tilde{\bm{X}}_t = \bm{X}_{T-t}$ es la solución del
problema~\eqref{eq:schrodinger_problem}.
\end{proposicion}

\begin{proof}
Por~\eqref{eq:kl_procesos}, minimizar $D_{\mathrm{KL}}(\mathbb{P}\|\mathbb{R})$
sobre $\Pi(p_T,p_0)$ equivale a encontrar el drift $b^*$ que minimiza la
norma cuadrática integrada $\mathbb{E}[\int_0^T\|b-f_{\mathbb{R}}\|^2/g^2\,dt]$
sujeto a que $\mathbb{P}$ tenga marginales $p_T$ y $p_0$.

Se introduce el lagrangiano con multiplicadores funcionales $\varphi(\bm{x})$
y $\psi(\bm{x})$ para las restricciones de marginalidad. La condición de
optimalidad de primer orden en el drift es:
\[
  b^*(\bm{x},t) = f_{\mathbb{R}}(\bm{x},t) + g(t)^2\nabla_{\bm{x}}\phi(\bm{x},t),
\]
donde $\phi$ satisface la ecuación de Kolmogorov hacia atrás respecto a
$\mathbb{R}$:
\[
  \partial_t\phi + f_{\mathbb{R}}\cdot\nabla\phi + \frac{g^2}{2}\Delta\phi = 0,
  \quad \phi(\bm{x},T) = \log\frac{p_0(\bm{x})}{p_T(\bm{x})}.
\]
Para $\mathbb{R}$ el forward OU corriendo hacia atrás (con $f_{\mathbb{R}} = \bm{x}$),
la función $\phi(\bm{x},t) = \log(p_{T-t}(\bm{x})/p_T(\bm{x}))$ satisface
esa ecuación (verificable por sustitución directa usando la Fokker-Planck del OU).
Luego $b^*(\bm{x},t) = \bm{x} + g^2\nabla\log p_{T-t}(\bm{x})$,
que es exactamente el drift del proceso de Anderson~\eqref{eq:reverse_ou}. $\square$
\end{proof}

\subsection{La dualidad de Sinkhorn}
\label{subsec:sinkhorn_dualidad}

El problema del puente de Schrödinger admite una formulación dual que produce
el \textbf{algoritmo de Sinkhorn}, que discretiza el problema en el espacio
de estados.

\subsubsection*{El problema primal discreto}

Dado el costo cuadrático $C_{ij} = \|\bm{x}_i-\bm{x}_j\|^2$, el kernel de
calor $K_{ij} = e^{-C_{ij}/\epsilon}$ (la discretización del kernel de
transición browniano con $\epsilon = 2g^2 T/n$ para $n$ pasos), y marginales
discretas $\bm{a} = (a_i)$ y $\bm{b} = (b_j)$, el problema de transporte
óptimo regularizado con entropía es:
\begin{equation}
  \bm{\pi}^* = \argmin_{\bm{\pi}\geq 0,\,\bm{\pi}\bm{1}=\bm{a},\,\bm{\pi}^\top\bm{1}=\bm{b}}
  \sum_{i,j}\pi_{ij}C_{ij} + \epsilon\sum_{i,j}\pi_{ij}\log\pi_{ij}.
  \label{eq:sinkhorn_primal}
\end{equation}

El término $\epsilon\sum_{ij}\pi_{ij}\log\pi_{ij} = \epsilon\,D_{\mathrm{KL}}(\bm{\pi}\|\bm{1})$
es la regularización entrópica: favorece planes de transporte difusos
(como el kernel browniano) sobre los concentrados en trayectorias únicas.

\subsubsection*{El problema dual: potenciales de Sinkhorn}

El dual del problema~\eqref{eq:sinkhorn_primal} es:
\begin{equation}
  \max_{\bm{f}\in\mathbb{R}^n,\,\bm{g}\in\mathbb{R}^n}
  \bm{f}\cdot\bm{a} + \bm{g}\cdot\bm{b}
  - \epsilon\sum_{i,j}e^{(f_i+g_j-C_{ij})/\epsilon},
  \label{eq:sinkhorn_dual}
\end{equation}
donde $\bm{f}$ y $\bm{g}$ son los \textbf{potenciales de Sinkhorn}.
Los potenciales tienen interpretación como multiplicadores de Lagrange de las
restricciones de marginalidad.

\begin{proposicion}{Forma de la solución primal y relación primal-dual}
\label{prop:sinkhorn_solucion}
La solución del problema~\eqref{eq:sinkhorn_primal} tiene la forma de
\textbf{Sinkhorn}:
\begin{equation}
  \pi_{ij}^* = u_i^*\,K_{ij}\,v_j^*,
  \label{eq:sinkhorn_forma}
\end{equation}
donde $\bm{u}^*$ y $\bm{v}^*$ son vectores positivos, relacionados con los
potenciales duales por $f_i = \epsilon\log u_i$ y $g_j = \epsilon\log v_j$.
Las condiciones de marginalidad $\bm{\pi}^*\bm{1} = \bm{a}$ y
$(\bm{\pi}^*)^\top\bm{1} = \bm{b}$ se traducen en:
$\bm{u}\odot(K\bm{v}) = \bm{a}$ y $\bm{v}\odot(K^\top\bm{u}) = \bm{b}$.
\end{proposicion}

\begin{proof}
Las condiciones KKT del primal~\eqref{eq:sinkhorn_primal} dan:
$\partial_{\pi_{ij}}(\text{Lagrangiano}) = C_{ij} + \epsilon(1+\log\pi_{ij})
- \lambda_i - \nu_j = 0$,
luego $\pi_{ij}^* = e^{(\lambda_i+\nu_j-C_{ij})/\epsilon - 1}
= e^{(\lambda_i-1/2)/\epsilon}K_{ij}e^{(\nu_j-1/2)/\epsilon}
=: u_i K_{ij} v_j$, con $u_i = e^{(\lambda_i-\epsilon/2)/\epsilon}$ y
análogamente para $v_j$. Las condiciones de marginalidad imponen
$u_i\sum_j K_{ij}v_j = a_i$ y $v_j\sum_i K_{ij}u_i = b_j$, que en notación
vectorial son $\bm{u}\odot(K\bm{v}) = \bm{a}$ y $\bm{v}\odot(K^\top\bm{u}) = \bm{b}$.
$\square$
\end{proof}

\subsubsection*{El algoritmo de Sinkhorn}

Las condiciones de marginalidad sugieren el siguiente método de punto fijo
alternado (\textbf{algoritmo de Sinkhorn}, Sinkhorn 1967, redescubierto por
Cuturi 2013):

\begin{algorithm}[H]
\caption{Algoritmo de Sinkhorn}
\label{alg:sinkhorn}
\KwIn{Marginales $\bm{a}$, $\bm{b}$; kernel $K_{ij} = e^{-C_{ij}/\epsilon}$;
      tolerancia $\delta$}
\KwOut{Plan óptimo $\bm{\pi}^* = \mathrm{diag}(\bm{u}^*)K\mathrm{diag}(\bm{v}^*)$}
$\bm{v}^{(0)} \leftarrow \bm{1}$\;
\For{$k = 0, 1, 2, \ldots$ hasta convergencia}{
  $\bm{u}^{(k+1)} \leftarrow \bm{a} \oslash (K\bm{v}^{(k)})$\;
  $\bm{v}^{(k+1)} \leftarrow \bm{b} \oslash (K^\top\bm{u}^{(k+1)})$\;
}
\Return $\bm{\pi}^* = \mathrm{diag}(\bm{u}^*)K\mathrm{diag}(\bm{v}^*)$\;
\end{algorithm}

donde $\oslash$ es la división elemento a elemento.

\begin{proposicion}{Convergencia del algoritmo de Sinkhorn}
\label{prop:sinkhorn_convergencia}
El Algoritmo~\ref{alg:sinkhorn} converge linealmente con factor de contracción
$\kappa = \tanh(\lambda_K/4)^2 < 1$, donde $\lambda_K = \log(\max_{ij}K_{ij}/
\min_{ij}K_{ij})$ es el contraste del kernel.
\end{proposicion}

\begin{proof}[Esquema]
Se trabaja en escala logarítmica $\bm{f}^{(k)} = \epsilon\log\bm{u}^{(k)}$,
$\bm{g}^{(k)} = \epsilon\log\bm{v}^{(k)}$. Cada iteración de Sinkhorn es una
proyección alternada en el espacio de potenciales duales sobre los conjuntos
afines $\{\bm{f}: \sum_j e^{(f_i+g_j-C_{ij})/\epsilon} = a_i\,\forall i\}$ y
análogo para $\bm{g}$. La contracción de la proyección alternada sobre dos
conjuntos convexos es la razón entre la norma dual del kernel y la norma
primal, que da la tasa $\kappa$. Para $\epsilon$ grande (kernel suave),
$\kappa$ es pequeño y la convergencia es rápida. $\square$
\end{proof}

\subsection{El límite de difusión cero: transporte óptimo de Monge}
\label{subsec:ot_limite}

\begin{teorema}{Convergencia al transporte óptimo de Monge}
\label{thm:ot_limite}
En el límite $\epsilon \to 0$, el plan óptimo regularizado $\bm{\pi}^*_\epsilon$
converge al plan de \textbf{transporte óptimo de Monge-Kantorovich} sin
regularización:
\begin{equation}
  \bm{\pi}^*_0 = \argmin_{\bm{\pi}\geq 0,\,\bm{\pi}\bm{1}=\bm{a},\,
  \bm{\pi}^\top\bm{1}=\bm{b}}
  \sum_{i,j}\pi_{ij}C_{ij},
  \label{eq:ot_sin_reg}
\end{equation}
en el sentido que $W_2(\bm{\pi}^*_\epsilon,\bm{\pi}^*_0)\to 0$ con
$W_2(\bm{\pi}^*_\epsilon,\bm{\pi}^*_0) = O(\epsilon\log(1/\epsilon))$.
\end{teorema}

\begin{proof}[Esquema]
Cuando $\epsilon \to 0$, el kernel $K_{ij} = e^{-C_{ij}/\epsilon}$ se concentra
en cero: $K_{ij} \to 0$ para $C_{ij} > 0$ y $K_{ii} = 1$ para $C_{ii} = 0$.
El plan de Sinkhorn $\pi_{ij}^* = u_i K_{ij} v_j$ concentra la masa en los
pares $(i,j)$ con $C_{ij}$ mínimo compatibles con las marginales: el mapa
determinista $j = T^*(i)$ que minimiza el costo cuadrático. La regularización
entrópica $\epsilon\sum_{ij}\pi_{ij}\log\pi_{ij}$ actúa como barrera que
desplaza la solución en $O(\epsilon\log(1/\epsilon))$ de la solución exacta.
$\square$
\end{proof}

Para la distribución de datos $p_0 = \mathcal{N}(\bm{\mu}_0, \bm{\Sigma}_0)$
y prior $p_T = \mathcal{N}(\bm{0},\bm{I})$, el mapa de transporte óptimo de
Monge bajo costo cuadrático es el mapa afín:
\begin{equation}
  T^*(\bm{x}) = \bm{\Sigma}_0^{-1/2}
  \left(\bm{\Sigma}_0^{1/2}\bm{I}\bm{\Sigma}_0^{1/2}\right)^{1/2}
  \bm{\Sigma}_0^{-1/2}(\bm{x}-\bm{\mu}_0)
  = \bm{\Sigma}_0^{-1/2}(\bm{x}-\bm{\mu}_0),
  \label{eq:monge_gaussiana}
\end{equation}
y las trayectorias de transporte óptimo son líneas rectas:
$\bm{x}(t) = (1-t)\bm{x}_0 + t T^*(\bm{x}_0)$,
con campo de velocidades $\dot{\bm{x}} = T^*(\bm{x}_0)-\bm{x}_0$ constante a
lo largo de cada trayectoria.

\subsection{Flow matching: aprendizaje del campo de velocidades}
\label{subsec:flow_matching}

La ODE de flujo de probabilidad~\eqref{eq:prob_flow_ode} y el límite de
transporte óptimo sugieren aprender directamente el campo vectorial $v(\bm{x},t)$
en lugar del score.

\begin{definicion}{Flow matching (Lipman et al., 2022; Liu et al., 2022)}
\label{def:flow_matching}
El \textbf{flow matching} aprende el campo $v_{\bm{\phi}}(\bm{x},t)$ de la ODE
$d\bm{x}/dt = v(\bm{x},t)$ minimizando:
\begin{equation}
  \mathcal{L}_{\mathrm{FM}}(\bm{\phi})
  = \mathbb{E}_{t,\bm{x}_0\sim p_0,\bm{x}_T\sim p_T}
  \left[\|v_{\bm{\phi}}(\bm{x}(t),t) - (\bm{x}_T-\bm{x}_0)\|^2\right],
  \label{eq:fm_loss}
\end{equation}
donde $\bm{x}(t) = (1-t)\bm{x}_0 + t\bm{x}_T$ es la interpolación lineal y
$\bm{x}_T-\bm{x}_0$ es la velocidad objetivo.
\end{definicion}

\begin{proposicion}{El campo de velocidades óptimo del flow matching}
\label{prop:fm_campo}
El minimizador de~\eqref{eq:fm_loss} es el campo de velocidades condicional:
\[
  v^*(\bm{x},t)
  = \mathbb{E}_{\bm{x}_0,\bm{x}_T|\bm{x}(t)=\bm{x}}[\bm{x}_T-\bm{x}_0].
\]
Para pareo independiente ($\bm{x}_0 \perp \bm{x}_T$), las trayectorias de
distintos pares se cruzan y $v^*$ es un promedio sobre ellas. Para pareo
óptimo $\bm{x}_T = T^*(\bm{x}_0)$ (OT flow matching, Tong et al., 2023),
las trayectorias no se cruzan y $v^*(\bm{x},t) = T^*(\bm{x}_0) - \bm{x}_0$
es constante a lo largo de cada trayectoria.
\end{proposicion}

\begin{proof}
Expandiendo $\mathcal{L}_{\mathrm{FM}} = \mathbb{E}\|v_{\bm{\phi}}(\bm{x}(t),t)
- (\bm{x}_T-\bm{x}_0)\|^2$ y tomando la derivada funcional respecto a
$v_{\bm{\phi}}(\bm{x},t)$:
$v^*(\bm{x},t) = \argmin_v \mathbb{E}[\|v - (\bm{x}_T-\bm{x}_0)\|^2 |
\bm{x}(t)=\bm{x}] = \mathbb{E}[\bm{x}_T-\bm{x}_0|\bm{x}(t)=\bm{x}]$.
Para OT flow matching, el pareo $(\bm{x}_0,T^*(\bm{x}_0))$ es determinista dado
$\bm{x}_0$, y las trayectorias $\{(1-t)\bm{x}_0+tT^*(\bm{x}_0):t\in[0,1]\}$
no se cruzan entre distintos $\bm{x}_0$ (propiedad del mapa de Monge). Luego
$\bm{x}(t)$ determina únicamente $\bm{x}_0$ y la esperanza colapsa al valor
exacto. $\square$
\end{proof}

\subsubsection*{Comparación cuantitativa: difusión vs.\ flow matching}

\begin{table}[h]
\centering
\begin{tabular}{lllll}
\toprule
\textbf{Método} & $\epsilon$ & \textbf{Trayectorias} & \textbf{Pasos} & \textbf{$v^*$ complejidad} \\
\midrule
DDPM & $\infty$ & Aleatorias & $1000$ & Alta \\
Score SDE & grande & Aleatorias & $10$-$100$ & Alta \\
Prob. flow ODE & grande & Deterministas & $10$-$100$ & Alta \\
OT flow matching & $0^+$ & Rectas & $1$-$10$ & Baja \\
\midrule
Schrödinger bridge & $\epsilon>0$ & Suavizadas & $10$-$100$ & Media \\
\bottomrule
\end{tabular}
\caption{Comparación de métodos generativos basados en transporte.}
\label{tab:comparacion}
\end{table}

El costo del campo de velocidades en OT flow matching es mínimo porque las
trayectorias rectas no se cruzan: la red no necesita resolver la ambigüedad
de qué modo pertenece cada punto de la trayectoria. Esto explica por qué los
modelos de tercera generación (Stable Diffusion 3, Flux) usan flow matching
en lugar de difusión estándar.

%───────────────────────────────────────────────────────────────────────────────
\section{Resumen del Capítulo 11}
\label{sec:resumen11}
%───────────────────────────────────────────────────────────────────────────────

El capítulo establece cuatro resultados encadenados. Primero, la ecuación de
Fokker-Planck derivada por dualidad $L^2$ cierra la deuda del Capítulo~\ref{ch:cap10} y
la Proposición~\ref{prop:gibbs_estacionaria} cierra la del Capítulo~\ref{ch:cap5}. La
identidad~\eqref{eq:identidad_difusiva} —que el término de segundo orden de la
Fokker-Planck es la mitad de la divergencia de $(\nabla\cdot\bm{D})p + \bm{D}\nabla p$—
es el instrumento algebraico central de todo el capítulo.

Segundo, el Teorema~\ref{thm:anderson_general} establece la SDE del proceso
reverso en su forma más general, con $\bm{\Sigma}(\bm{x},t)$ arbitrario. El
término $\nabla\cdot\bm{D}$ es la corrección por heterogeneidad espacial de la
difusión: cero para difusión isótropa, no nulo para medios heterogéneos (física
de reactores, medios porosos). Se presenta en dos rutas: la Ruta A (verificación
de Fokker-Planck) es elemental; la Ruta B (Girsanov) revela que el score es la
derivada de Radon-Nikodym que corrige el browniano reverso para tener las
marginales correctas.

Tercero, los Teoremas~\ref{thm:hyvarinen} y~\ref{thm:dsm} convierten el score
en la pérdida de predicción de ruido~\eqref{eq:ddpm_loss}: computable, escalable,
y exactamente DDPM. La discusión sobre correlación aclara que la difusión
isótropa no destruye la correlación de los datos sino que la traslada al score,
que la red aprende implícitamente.

Cuarto, la Proposición~\ref{prop:anderson_schrodinger} revela que el proceso de
Anderson es el puente de Schrödinger: el proceso más cercano al forward en KL
entre leyes de procesos. La dualidad de Sinkhorn (Algoritmo~\ref{alg:sinkhorn})
discretiza ese problema variacional y converge al mapa de transporte óptimo de
Monge en el límite $\epsilon\to 0$ (Teorema~\ref{thm:ot_limite}). El flow
matching aprende ese mapa directamente con trayectorias rectas y sin cruce,
reduciendo el número de pasos de muestreo de $10^3$ (DDPM) a $O(10)$.

La conexión hacia el \textbf{Capítulo~\ref{ch:cap12}} es ahora completa: forward process
(proceso OU, Capítulo~\ref{ch:cap10}), reverse SDE (Teorema~\ref{thm:anderson_general}),
pérdida de entrenamiento (Teorema~\ref{thm:dsm}), arquitectura (U-Net/transformer
equivariante, Capítulo~\ref{ch:cap9}), y el marco unificador del puente de Schrödinger que
interpola entre difusión (DDPM) y transporte (flow matching).

%───────────────────────────────────────────────────────────────────────────────
\label{sec:estrella11}
\begin{estrella}{El reverse SDE en reactores heterogéneos y modelos generativos de flujo neutrónico}

\textbf{Forward process como modelo de difusión de neutrones perturbada.}
El flujo neutrónico perturbado $\delta\Phi(\bm{r},t)$ en torno al estado
estacionario del IAN-R1 satisface, en la aproximación de un grupo y geometría
cilíndrica, una ecuación de difusión análoga a la Fokker-Planck con difusividad
$D(\bm{r})$ variable (zona de combustible, reflector de berilio, estructura):
\[
  \frac{\partial p_t}{\partial t}
  = -\nabla\cdot\!\left[(\nabla D)\,p_t + D(\bm{r})\nabla p_t\right]
  + \Sigma_a\,p_t = \mathcal{A}^*p_t + \Sigma_a\,p_t,
\]
donde $p_t(\bm{r})$ es la densidad de probabilidad de encontrar un neutrón en
la posición $\bm{r}$ en el tiempo $t$ (proporcional al flujo local). El término
$\nabla D$ es exactamente la contribución $\nabla\cdot\bm{D}$ del Teorema
general~\ref{thm:anderson_general}: no puede ignorarse en las interfaces
combustible-moderador-reflector.\medskip

\textbf{El reverse SDE como modelo generativo de distribuciones de flujo.}
La SDE~\eqref{eq:reverse_general} con el tensor de difusión $D(\bm{r})$ del
reactor proporciona un método novedoso para generar distribuciones de flujo
físicamente consistentes para distintas configuraciones de barras de control,
sin ejecutar el código de difusión completo. El proceso forward (perturbación
del flujo desde la distribución nominal hacia el ruido) y el reverse (generación
desde el ruido hacia una distribución de flujo física) forman un par análogo al
forward/reverse de los modelos de difusión de imágenes, pero sobre el espacio
de distribuciones de neutrones. El score $\nabla\log p_t(\bm{r})$ aprendido
por la red es el campo de fuerza que corrige la difusión browniana isótropa para
respetar la geometría heterogénea del reactor.\medskip

\textbf{Transporte óptimo entre distribuciones de flujo.}
El flow matching del Capítulo~\ref{ch:cap12} (OT flow matching) puede
aplicarse para transportar óptimamente entre distribuciones de flujo
correspondientes a distintas condiciones de operación: de flujo a potencia
$P_1$ a flujo a potencia $P_2$, o de configuración de barras $\mathcal{B}_1$
a $\mathcal{B}_2$. El plan de transporte óptimo discretizado con el Algoritmo
de Sinkhorn~\ref{alg:sinkhorn} sobre una malla del reactor (con costo
cuadrático $C_{ij} = \|\bm{r}_i - \bm{r}_j\|^2$) produce el campo de
velocidades de menor energía que conecta ambas distribuciones: el camino
más corto en el espacio de distribuciones de flujo, con interpretación directa
como la manobra de barras de control más eficiente para transicionar entre
los dos estados.

El score matching para el IAN-R1 tiene una ventaja particular respecto a
otros reactores: el conjunto de simulaciones de referencia (MCNP o difusión
de dos grupos) es pequeño ($N \sim 10^3$--$10^4$), pero el denoising score
matching~\eqref{eq:ddpm_loss} es un estimador eficiente incluso con pocos datos
porque aprovecha la regularidad suave de las distribuciones de flujo (límites
continuos, simetría cilíndrica, decaimiento en el reflector).
\end{estrella}

%───────────────────────────────────────────────────────────────────────────────
\section*{Ejercicios computacionales}
\addcontentsline{toc}{section}{Ejercicios computacionales}
%───────────────────────────────────────────────────────────────────────────────

\begin{ejerciciocomp}{Verificación del teorema de Anderson por simulación directa}
\textbf{Objetivo:} verificar numéricamente que el proceso reverso~\eqref{eq:reverse_ou}
para el OU reproduce $p_0$ a partir de muestras de $p_T$, usando el score exacto
gaussiano.

\medskip\noindent\textbf{Algoritmo:}
\begin{enumerate}
  \item Tomar $p_0 = \frac{1}{2}\mathcal{N}(-2, 0.5^2) + \frac{1}{2}\mathcal{N}(2, 0.5^2)$
        (bimodal). Simular el forward process OU ($\alpha=1$, $\sigma=\sqrt{2}$)
        con Euler-Maruyama desde $t=0$ hasta $T=4$ usando $\Delta t = 0.01$.
        Verificar que la distribución en $t=T$ es aproximadamente $\mathcal{N}(0,1)$.
  \item En cada tiempo $t$, la distribución marginal del forward OU es una
        mezcla gaussiana con parámetros conocidos analíticamente (del Ejercicio
        Resuelto~\ref{ej:ou_distribucion}). Calcular el score exacto
        $\nabla\log p_t(x) = \sum_k w_k(x,t)\cdot(-(x-\mu_k(t))/\nu_k(t))$
        como promedio ponderado de las componentes.
  \item Simular el reverse SDE~\eqref{eq:reverse_ou} con el score exacto,
        desde $t=0$ ($\tilde{p}_0 = p_T \approx \mathcal{N}(0,1)$) hasta $t=T$
        ($\tilde{p}_T$ debe aproximar $p_0$). Usar $N=10^3$ partículas y
        $\Delta t = 0.01$.
  \item Comparar el histograma de $\tilde{X}_T$ con $p_0$ usando la divergencia
        KL empírica. ¿El reverse SDE recupera la bimodalidad? ¿Cuánto influye
        el paso $\Delta t$?
\end{enumerate}
\medskip\noindent\textbf{Verificación:} $D_\mathrm{KL}(\hat{p}_0\|p_0) < 0.05$
con $N = 10^3$ partículas y $\Delta t = 0.01$.
\end{ejerciciocomp}

\begin{ejerciciocomp}{Denoising score matching: entrenamiento y muestreo}
\textbf{Objetivo:} entrenar un score network con la pérdida~\eqref{eq:ddpm_loss}
y usarlo en el reverse SDE para generar muestras de $p_0$.

\medskip\noindent\textbf{Algoritmo:}
\begin{enumerate}
  \item Usar la misma $p_0$ bimodal del ejercicio anterior. Entrenar un MLP
        $(2, 64, 64, 1)$ con entradas $(x, \sigma)$ y salida $\varepsilon_\phi(x,\sigma)$
        usando la pérdida $\mathbb{E}_{p_0,\varepsilon}[\|\varepsilon_\phi(x +
        \sigma\varepsilon, \sigma) - \varepsilon\|^2]$ sobre un schedule de
        $\sigma \in [0.01, 3.0]$ con $10^4$ iteraciones de Adam.
  \item Convertir el score estimado: $\hat{s}(x,\sigma) = -\varepsilon_\phi(x,\sigma)/\sigma$.
        Comparar $\hat{s}(x, \sigma=0.3)$ con el score exacto gaussiano sobre
        la malla $x \in [-5, 5]$.
  \item Muestrear 500 puntos con el reverse SDE de Euler-Maruyama (200 pasos)
        y comparar con $p_0$ mediante KDE. ¿Reproduce las dos modas?
  \item Variar el número de pasos de integración en $\{20, 50, 100, 200\}$.
        ¿Cómo afecta la calidad de las muestras?
\end{enumerate}
\end{ejerciciocomp}

%───────────────────────────────────────────────────────────────────────────────
\begin{notasbib}
La ecuación de Fokker-Planck fue introducida por Fokker (1914) y Planck (1917)
en el contexto de la difusión browniana. La formulación moderna como adjunto
$L^2$ del generador es estándar; véase Risken (1989)~\cite{risken1989fokker}
para una presentación completa orientada a física estadística. El H-teorema y la
estructura de producción de entropía relacionada con la información de Fisher
se discuten en Villani (2009)~\cite{villani2009optimal}.

El teorema de Anderson (Teorema~\ref{thm:anderson_general}) fue demostrado por
Anderson (1982)~\cite{anderson1982reverse} mediante la verificación directa
(Ruta A). La forma general con $\bm{D}(\bm{x},t)$ dependiente del estado no
siempre se presenta en los tutoriales de modelos de difusión; la exposición aquí
sigue a Song et al.\ (2021)~\cite{song2021score} y Särkkä y Solin
(2019)~\cite{sarkka2019applied}. La Ruta B vía el teorema de Girsanov
(Teorema~\ref{thm:girsanov}) sigue a la presentación de Föllmer y Wakolbinger
(1986) y Lehec (2013); véase también Tzen y Raginsky (2019) para la conexión
con la inferencia variacional.

El score matching implícito (Teorema~\ref{thm:hyvarinen}) fue introducido por
Hyvärinen (2005)~\cite{hyvarinen2005estimation}. El denoising score matching
(Teorema~\ref{thm:dsm}) fue introducido por Vincent (2011)~\cite{vincent2011connection}.
La pérdida de predicción de ruido~\eqref{eq:ddpm_loss} y el modelo DDPM son de
Ho et al.\ (2020)~\cite{ho2020denoising}. La formulación unificada mediante SDEs
(tanto el forward process como el reverse, incluyendo la ODE de flujo de
probabilidad) es de Song et al.\ (2021)~\cite{song2021score}.

El puente de Schrödinger fue formulado por Schrödinger (1931--1932) como un
problema de la mecánica cuántica en tiempo imaginario. Su formulación moderna en
términos de KL entre medidas de proceso es de Föllmer (1988)~\cite{follmer1988random}
y Leonard (2014)~\cite{leonard2014survey}. El algoritmo de Sinkhorn fue
introducido por Sinkhorn (1964)~\cite{sinkhorn1964relationship} para matrices
de doble estocástica y redescubierto en el contexto de transporte óptimo
regularizado por Cuturi (2013)~\cite{cuturi2013sinkhorn}. El flow matching fue
introducido simultáneamente por Lipman et al.\ (2022)~\cite{lipman2022flow} y
Liu et al.\ (2022)~\cite{liu2022flow}; el OT flow matching con pareo óptimo es
de Tong et al.\ (2023)~\cite{tong2023improving}.

La conexión entre la ecuación de difusión de neutrones y la Fokker-Planck en
medios heterogéneos (relevante para la Sección~\ref{sec:estrella11}) se discute
en el contexto de la ecuación de transporte de Boltzmann en Bell y Glasstone
(1970)~\cite{bell1970nuclear}. El término de difusioforesis en la ecuación de
transporte de partículas es tratado en Anderson (1989)~\cite{anderson1989colloid}.
\end{notasbib}

%───────────────────────────────────────────────────────────────────────────────
\section*{Ejercicios resueltos}
\addcontentsline{toc}{section}{Ejercicios resueltos}
%───────────────────────────────────────────────────────────────────────────────

\begin{ejercicio}
\label{ej:fp_ou}
\textbf{(Fokker-Planck para el proceso OU: familia gaussiana).}
Para $dX_t = -\alpha X_t\,dt + \sigma\,dW_t$, verificar que la familia
$p_t = \mathcal{N}(\mu_t,\nu_t)$ satisface la Fokker-Planck y encontrar
las ODEs para $\mu_t$, $\nu_t$.

\medskip\noindent\textit{Solución.}
La Fokker-Planck es $\partial_t p = \partial_x(\alpha x p) + \frac{\sigma^2}{2}
\partial_{xx}p$. Para $p = \mathcal{N}(\mu_t,\nu_t)$, identificando coeficientes
de $(x-\mu_t)^0$ y $(x-\mu_t)^2$:
$\dot{\mu}_t = -\alpha\mu_t$, $\dot{\nu}_t = -2\alpha\nu_t + \sigma^2$.
Soluciones: $\mu_t = \mu_0 e^{-\alpha t}$, $\nu_t = \nu_0 e^{-2\alpha t} +
\frac{\sigma^2}{2\alpha}(1-e^{-2\alpha t})$, en perfecta concordancia con la
Sección~\ref{subsec:ou}. $\square$
\end{ejercicio}

\begin{ejercicio}
\label{ej:anderson_diag}
\textbf{(Teorema de Anderson para difusión diagonal: aplicación en reactor).}

Considerar un sistema neutrónico 1D con difusividad variable $D(x) > 0$:
\[
  dX_t = v(X_t)\,dt + \sqrt{2D(X_t)}\,dW_t.
\]
\textbf{(a)} Calcular $\nabla\cdot\bm{D}$ para este caso escalar.
\textbf{(b)} Escribir el reverse SDE de Anderson.
\textbf{(c)} Mostrar que si el proceso satisface balance detallado con
distribución estacionaria $p^*$, entonces $v(x) = D'(x) + D(x)(\log p^*)'(x)$.

\medskip\noindent\textit{Solución.}
\textbf{(a)} $D(x,t) = 2D(x)$ (escalar), luego $\nabla\cdot\bm{D} = 2D'(x)$.

\textbf{(b)} Del Teorema~\ref{thm:anderson_general}:
$d\tilde{X}_t = [-v(\tilde{X}_t,T-t) + 2D'(\tilde{X}_t)
+ 2D(\tilde{X}_t)\partial_x\log p_{T-t}(\tilde{X}_t)]\,dt
+ \sqrt{2D(\tilde{X}_t)}\,d\tilde{W}_t$.

\textbf{(c)} Del balance detallado~\eqref{eq:balance_general} con $D_{11} = 2D(x)$:
$v(x) = \frac{1}{2}\cdot 2D'(x) + \frac{1}{2}\cdot 2D(x)(\log p^*)'(x)
= D'(x) + D(x)(\log p^*)'(x)$.
Este es el drift de la difusión de Smoluchowski en un potencial efectivo
$U(x) = -\log p^*(x)$, con el término adicional $D'(x)$ que no aparece si
$D$ es constante. En la física de reactores, $D'(x)$ representa la fuerza
difusioforética en interfaces entre materiales de distinta difusividad. $\square$
\end{ejercicio}

%───────────────────────────────────────────────────────────────────────────────
\section*{Ejercicios propuestos}
\addcontentsline{toc}{section}{Ejercicios propuestos}
%───────────────────────────────────────────────────────────────────────────────

\begin{ejercicio}
\textbf{(Reverse SDE para difusión con tensor de difusión completo).}
Sea $d\bm{X}_t = -\bm{A}\bm{X}_t\,dt + \bm{B}\,d\bm{W}_t$ con $\bm{A}$
y $\bm{B}$ matrices constantes, $\bm{D} = \bm{B}\bm{B}^\top$.
\begin{enumerate}[label=(\alph*)]
  \item Calcular $\nabla\cdot\bm{D}$ para este caso.
  \item Calcular la distribución marginal gaussiana $p_t =
        \mathcal{N}(\bm{\mu}_t,\bm{\Sigma}_t)$ resolviendo las ecuaciones
        de Lyapunov $\dot{\bm{\mu}}_t = -\bm{A}\bm{\mu}_t$ y
        $\dot{\bm{\Sigma}}_t = -\bm{A}\bm{\Sigma}_t - \bm{\Sigma}_t\bm{A}^\top + \bm{D}$.
  \item Escribir el reverse SDE explícito. ¿Qué forma tiene el score para
        $p_t$ gaussiana?
  \item Para $\bm{A} = \mathrm{diag}(1,2)$ y $\bm{D} = \mathrm{diag}(2,4)$,
        verificar que la distribución estacionaria es $\mathcal{N}(\bm{0},\bm{I})$
        y que el reverse SDE lleva $p_T = \mathcal{N}(\bm{0},\bm{I})$ de vuelta
        a $p_0$.
\end{enumerate}
\end{ejercicio}

\begin{ejercicio}
\textbf{(Sinkhorn y transporte óptimo entre gaussianas).}
Sean $p_0 = \mathcal{N}(-2,0.1)$ y $p_T = \mathcal{N}(2,0.1)$ discretizadas en
$n=50$ puntos en $[-4,4]$.
\begin{enumerate}[label=(\alph*)]
  \item Implementar el Algoritmo~\ref{alg:sinkhorn} para $\epsilon \in
        \{0.01, 0.1, 1.0\}$. Verificar que el costo $\sum_{ij}\pi_{ij}^*C_{ij}$
        converge a $16 = (-2-2)^2$ cuando $\epsilon\to 0$.
  \item Calcular el mapa de transporte inducido $T_\epsilon(x_i) =
        \sum_j \pi_{ij}^* x_j / a_i$ y compararlo con $T^*(x) = x+4$.
  \item Verificar numéricamente que el error $\|T_\epsilon - T^*\|_\infty =
        O(\epsilon\log(1/\epsilon))$ como predice el Teorema~\ref{thm:ot_limite}.
  \item Trazar la interpolación $\bm{x}(t) = (1-t)x_0 + tT^*(x_0)$ para 20
        valores de $x_0$ muestreados de $p_0$. ¿Se cruzan las trayectorias?
\end{enumerate}
\end{ejercicio}

\begin{ejercicio}
\textbf{(Flow matching vs.\ score matching para la misma tarea).}
Sea $p_0 = \frac{1}{2}\mathcal{N}(-3,0.25) + \frac{1}{2}\mathcal{N}(3,0.25)$
(mezcla bimodal) y $p_T = \mathcal{N}(0,1)$.
\begin{enumerate}[label=(\alph*)]
  \item Entrenar un MLP con la pérdida de denoising score matching~\eqref{eq:ddpm_loss}
        con $\sigma \in \{0.1, 0.5, 1.0\}$ y $10^4$ muestras.
  \item Entrenar un MLP con la pérdida de flow matching~\eqref{eq:fm_loss}
        con pareo independiente.
  \item Entrenar el mismo MLP con pareo OT (usando Sinkhorn del ejercicio
        anterior). Comparar la varianza del campo $v^*$ en los tres casos.
  \item Muestrear $10^3$ puntos de cada modelo integrando el reverse SDE
        (score matching) o la ODE (flow matching) con el método de Euler.
        ¿Cuántos pasos necesita cada método para producir muestras de calidad
        comparable?
\end{enumerate}
\end{ejercicio}

\begin{ejercicio}
\textbf{(Difusividad variable en un reactor esférico homogéneo).}
Considerar la ecuación de difusión de neutrones en 1D esférico con difusividad
$D(r) = D_0(1-r^2/R^2)$ (máxima en el centro, cero en el borde):
$dX_t = [D'(X_t) + D(X_t)\partial_x\log p^*(X_t)]\,dt + \sqrt{2D(X_t)}\,dW_t$.
\begin{enumerate}[label=(\alph*)]
  \item Verificar que la distribución estacionaria de esta SDE satisface
        $p^*(x) = Ce^{-\int_0^x v(x')/D(x')\,dx'}$ donde $v(x) = D'(x)$
        (sin potencial externo) es la fuerza de difusioforesis sola.
  \item Calcular $p^*(x)$ explícitamente para el perfil $D(r) = D_0(1-r^2/R^2)$.
  \item Escribir el reverse SDE para este sistema y verificar que lleva $p_T$
        de vuelta a $p^*(x)$.
  \item Simular el forward process (dispersión de neutrones en el reactor) y
        el reverse process con el score exacto. ¿Cuánto tiempo necesita el
        forward para llegar a $\mathcal{N}(0,\sigma^2)$?
\end{enumerate}
\end{ejercicio}
